\note{sa page 9}

Alors on veut que
\[ L^{n-i+\alpha}\, q_{i,\alpha} \colon M^{i} \longrightarrow M^{2n-i+2\alpha} \]
soit nul sur les $L^{\alpha'} P^{i-2\alpha'}$ pour $\alpha' \neq \alpha$, et que
sur $L^{\alpha} P^{i-2\alpha}$ il induise \ill{} le composé avec l'isomorphisme
\[ L^{n-i+2\alpha} \colon P^{i-2\alpha} \longrightarrow L^{n-i+2\alpha} P^{i-2\alpha} \subset M^{2n-i+2\alpha} . \]
\ill{} il s'ensuit aussitôt \ill{} que

\[ (6.9) \qquad \pi_{i,\alpha} = L^{\alpha}\, w_{i-2\alpha}\, L^{n-i+\alpha}\, q_{i,\alpha} = L^{\alpha}\, w_{i-2\alpha}\, L^{n-i+\alpha}\Bigl( 1 - \sum_{\alpha' > \alpha} \pi_{i,\alpha'} \Bigr) \]

ce qui donne bien $\pi_{i,\alpha} \in \mathcal{E}$. \struck{Ceci prouve (iv).}
Si $i < n$, il faut de même $\pi_{2n-i,\alpha} \in \mathcal{E}$. \ill{} On a
$\pi_{2n-i,\alpha} \colon M^{2n-i} \to M^{2n-i+2\alpha}$ \ill{} de la façon

\[ (6.10) \qquad \pi_{2n-i,\alpha} = L^{n-i}\, \pi_{i,\alpha}\, w_{i} \ \struck{\ill{}} \]

\emph{Reste} : prouver (iii) $\Lambda \in \mathcal{E}$. Comme on a

\[ (6.11) \qquad \Lambda = \sum_{j} \struck{\ill{}}\ \Lambda\, \pi_{j,\alpha}\ \struck{\ill{}} \]

il suffit de prouver que $\Lambda \pi_{i,\alpha} \struck{\ill{}} \in \mathcal{E}$.

